% !TEX root = ../main.tex
\chapter{Analiza tematu}
\label{ch:analiza}

% Zgodnie z wytycznymi szablonu:
% - sformułowanie problemu (state of the art)
% - poszerzone studia literaturowe
% - opis znanych rozwiązań i algorytmów
% - osadzenie pracy w kontekście
% Rozdział powinien być wysycony cytowaniami.

W tym rozdziale przedstawiono tło problemu, najważniejsze typy nadużyć oraz
dotychczasowe podejścia stosowane do ich wykrywania w komunikacji tekstowej.
Omówienie porządkuje terminologię i wskazuje ograniczenia metod regułowych
oraz modeli opartych na uczeniu maszynowym.


\section{Sformułowanie problemu}
\label{sec:problem}

Komunikacja tekstowa w aplikacjach typu C2C (\english{consumer-to-consumer}) jest
kluczowym elementem finalizacji transakcji, ale jednocześnie stanowi kanał podatny
na nadużycia. Użytkownicy próbują przenosić rozmowę i płatność poza platformę,
wyłudzać dane lub środki finansowe oraz publikować treści toksyczne. Zjawiska te
obniżają bezpieczeństwo, zaufanie i~retencję użytkowników, a~po stronie operatora
prowadzą do strat finansowych i~wzrostu kosztów moderacji.

Problem badawczy można sformułować jako znalezienie najlepszego podejścia do automatycznego wykrywania nadużyć
w~polskojęzycznych wiadomościach czatowych z~jednoczesnym zachowaniem wysokiej
skuteczności i~niskiej liczby fałszywych alarmów. W~praktyce oznacza to potrzebę
rozróżniania co najmniej czterech klas ryzykownych zachowań:

\begin{itemize}
\item \textbf{Obejście platformy} (\english{platform bypass}): próby przeniesienia komunikacji i~płatności poza platformę, np. poprzez podawanie numerów telefonu, adresów e-mail, linków do zewnętrznych serwisów płatniczych lub komunikatorów.
\item \textbf{Oszustwa} (\english{fraud}): próby wyłudzenia danych osobowych, informacji o~karcie płatniczej, danych do logowania.
\item \textbf{Toksyczność} (\english{toxicity}): używanie wulgaryzmów, obraźliwych zwrotów, mowy nienawiści lub innych form agresywnej komunikacji, które mogą zniechęcać użytkowników do korzystania z~platformy.
\item \textbf{Inne}: np. spam, próby manipulacji opiniami, publikowanie treści niezgodnych z~regulaminem.

\end{itemize}

W~ramach pracy problem ten jest analizowany jako zadanie klasyfikacji tekstu, dla
którego porównywane są podejścia regułowe i~modele uczenia maszynowego. Kluczowym
wyzwaniem pozostaje poprawna identyfikacja intencji nadawcy mimo skrótów,
kolokwializmów, błędów językowych i~celowego maskowania treści.
% Trzy główne kategorie nadużyć (z prezentacji):
% - obejście platformy (platform bypass)
% - fraud (BLIK, wyłudzenia)
% - toksyczność (wulgaryzmy, mowa nienawiści)

\section{Istniejące rozwiązania w~moderacji treści platform internetowych}
\label{sec:sota}

Detekcja nadużyć tekstowych w~systemach internetowych jest aktywnie badanym i~wdrażanym
problemem na~świecie. Stosowane rozwiązania można sklasyfikować według podejścia
technicznego: od~filtrowania regułowego, przez narzędzia oparte na~uczeniu maszynowym,
po~systemy hybrydowe.

\paragraph{Filtrowanie oparte na~regułach i~słownikach.}
Najprostszym i~najstarszym podejściem jest utrzymywanie słowników zabronionych słów
i~wyrażeń regularnych wykrywających dane kontaktowe. Platformy C2C stosują tego
rodzaju listy blokujące jako pierwszą linię obrony. Zaletą jest deterministyczność
i~brak konieczności danych treningowych; wadą jest niezdolność do~wykrywania
obejść słownikowych (celowe literówki, niestandardowe spacje w~numerach telefonów,
zastępowanie znaków), wysoka liczba fałszywych alarmów oraz konieczność
nieustannej ręcznej aktualizacji.

\paragraph{Uczenie maszynowe i~modele transformerowe.}
Dynamiczny rozwój dziedziny NLP przyniósł nowe rozwiązania.
Google i~Jigsaw od~2017~roku rozwijają \textit{Perspective API}
\cite{bib:perspective}: usługę oceny toksyczności tekstu opartą na~modelach
sieci neuronowych, dostępną przez REST API i~używaną przez tysiące portali
do~filtrowania komentarzy. Wulczyn, Thain i~Dixon \cite{bib:wulczyn2017} wykazali,
że~modele uczenia maszynowego trenowane na~dużych zbiorach komentarzy osiągają
skuteczność porównywalną z~ludzką oceną w~wykrywaniu ataków personalnych.
Duże platformy cyfrowe coraz powszechniej stosują wielojęzyczne modele
do~automatycznej moderacji treści, łącząc klasyfikację automatyczną
z~moderacją ludzką dla~przypadków granicznych.
\cite{bib:perspective,bib:wulczyn2017}.

\paragraph{Duże modele językowe jako klasyfikatory.}
Rozwój dużych modeli językowych (\english{Large Language Models}, LLM) otworzył
kolejne podejście do~moderacji treści: zamiast trenować dedykowany klasyfikator,
wiadomość przesyłana jest przez API do~modelu ogólnego wraz z~instrukcją w~języku
naturalnym opisującą, jakie treści należy uznać za~nadużycie.
OpenAI udostępnia dedykowany punkt końcowy \textit{Moderation API}
\cite{bib:openai-moderation}, który klasyfikuje tekst w~kilkunastu kategoriach
zagrożeń (mowa nienawiści, przemoc, treści seksualne i~inne) bez konieczności
treningu i~bez opłat za~wywołanie.
Badania wykazały, że~modele GPT stosowane w~trybie zero-jedynkowym dorównują
lub~przewyższają adnotatorów ludzkich w~zadaniach klasyfikacji tekstu
\cite{bib:gilardi2023}, co~czyni je~atrakcyjną alternatywą wszędzie tam,
gdzie dostępność danych treningowych jest ograniczona.

Podejście to~ma jednak istotne ograniczenia w~kontekście produkcyjnym.
Po~pierwsze, każde wywołanie API wiąże się z~opóźnieniem rzędu setek
milisekund i~kosztem proporcjonalnym do~liczby tokenów, co~przy setkach
tysięcy wiadomości dziennie generuje znaczące koszty operacyjne.
Po~drugie, wyniki są~niedeterministyczne: ta~sama wiadomość może zostać
oceniona inaczej przy~kolejnym wywołaniu. Po~trzecie, i~najważniejsze
z~perspektywy RODO, treści wiadomości użytkowników musiałyby być przesyłane
na~serwery zewnętrznego dostawcy poza Europejskim Obszarem Gospodarczym.
Po~czwarte, gotowe modele ogólne nie~są~dostrojone do~specyfiki
polskojęzycznego czatu marketplace, co~ogranicza ich skuteczność
dla~wzorców obejścia platformy nieobecnych w~danych treningowych.
Z~tych powodów podejście LLM-as-classifier nie~zostało uwzględnione
jako moduł produkcyjny w~niniejszej pracy; szczegółowa analiza ograniczeń
zawarta jest w~sekcji~\ref{sec:llm-ograniczenia}.

\paragraph{Komercyjne narzędzia do~wykrywania oszustw.}
W~obszarze \english{fraud detection} funkcjonuje rozwinięty rynek komercyjnych
rozwiązań (m.in.\ Sift, Kount, Signifyd), które oceniają ryzyko transakcji
i~kont użytkowników na~podstawie sygnałów behawioralnych i~kontekstowych.
Podejścia te są zorientowane na~systemy płatnicze i~wykrywanie kont syntetycznych,
nie zaś na~analizę treści wiadomości czatowych, co~wyraźnie odróżnia je
od~celu niniejszej pracy.

\paragraph{Specyfika platform C2C i~luka dla~języka polskiego.}
Mimo bogactwa gotowych rozwiązań dla~języka angielskiego, polskojęzyczny kontekst
czatu C2C pozostaje słabo pokryty. Dostępne polskie zasoby NLP,
m.in.\ zbiory z~konkursów PolEval dotyczące mowy nienawiści,
obejmują głównie dane z~mediów społecznościowych, nie zaś specyficzny
styl komunikacji kupujący-sprzedający. Gotowe modele wielojęzyczne
nie~odzwierciedlają polskich kolokwializmów ani~specyficznych wzorców
obejścia prowizji platform marketplace. Stąd wynika potrzeba stworzenia
własnego systemu i~własnego zbioru danych, co~stanowi jedno z~kluczowych
uzasadnień niniejszej pracy.

\section{Duże modele językowe jako narzędzia moderacji treści: analiza ograniczeń}
\label{sec:llm-ograniczenia}

W podrozdziale \ref{sec:sota} opisano podejście \english{LLM-as-classifier} jako jeden
ze~współczesnych trendów w~moderacji treści: wiadomość przesyłana jest do~zewnętrznego
API modelu GPT-4 (OpenAI), Claude (Anthropic) lub~Gemini (Google DeepMind) wraz
z~instrukcją opisującą politykę regulaminu, a~model zwraca decyzję moderacyjną.
Podejście jest atrakcyjne w~fazie prototypowania, ponieważ eliminuje konieczność
gromadzenia danych treningowych i~pozwala w~kilka godzin wdrożyć funkcjonujący
system \cite{bib:gilardi2023}. Poniżej zebrano szczegółową analizę ograniczeń,
które wykluczają to~podejście z~użycia produkcyjnego w~kontekście polskojęzycznego
czatu marketplace.

\paragraph{Ograniczenie prawne: RODO i~suwerenność danych.}
Każde wywołanie API moderacji wymaga transmisji treści wiadomości do~serwerów
zewnętrznego dostawcy. Dla~usług OpenAI, Anthropic i~Google Gemini serwery
podstawowe zlokalizowane są~poza Europejskim Obszarem Gospodarczym~(EOG).
Rozporządzenie RODO w~art.~46 uzależnia transfer danych osobowych do~krajów
trzecich od~spełnienia dodatkowych wymogów (standardowe klauzule umowne,
wiążące reguły korporacyjne lub~decyzja stwierdzająca odpowiedni stopień
ochrony). Treści wiadomości czatowych mogą zawierać dane osobowe
w~rozumieniu RODO: imiona, numery telefonów, adresy e-mail, a~nawet
dane płatnicze w~przypadku wiadomości stanowiących próbę wyłudzenia.
Użytkownicy platformy nie~wyrażają przy~rejestracji zgody na~przekazywanie
ich danych do~modeli AI zewnętrznych dostawców, a~retroaktywne uzyskanie
takiej zgody jest zarówno trudne technicznie, jak i~wątpliwe z~perspektywy
zasady minimalizacji danych. Lokalne wdrożenie modeli eliminuje to~ryzyko
w~całości: żadne dane użytkowników nie~opuszczają infrastruktury operatora.

\paragraph{Koszt i~skalowalność.}
Komercyjne modele LLM rozliczane są~w~modelu \english{pay-per-token}.
Przy~założeniu średniej długości wiadomości czatowej rzędu 20~tokenów
i~kosztu modelu ekonomicznego klasy GPT-3.5 (ok.~\$0{,}002/1\,000~tokenów
wejścia), milion wiadomości dziennie generuje koszt rzędu 40~\$~na~dobę;
przy~modelu GPT-4o, kilkukrotnie więcej. Dla~platform z~ruchem
rzędu dziesiątek milionów wiadomości miesięcznie wydatki przekraczają
tysiące dolarów i~rosną liniowo z~ruchem. Własny model wdrożony lokalnie
ponosi jednorazowy koszt infrastruktury niezależny od~liczby wywołań,
co~przy~dużym ruchu oznacza wielokrotnie niższy.

\paragraph{Opóźnienie i~uzależnienie od~dostępności zewnętrznego API.}
Wywołania zewnętrznych API LLM charakteryzują się opóźnieniami rzędu
200-2000~ms, zależnie od~obciążenia serwerów dostawcy i~długości
generowanej odpowiedzi. Dla~aplikacji czatowej oczekujący użytkownik
odczuwa każde opóźnienie powyżej kilkuset milisekund jako zakłócenie
płynności rozmowy. Co~więcej, system moderacji staje się zależny od~ciągłości
działania zewnętrznego API: awaria, przekroczenie limitu wywołań
(\english{rate limit}) lub~planowana przerwa techniczna dostawcy
mogą całkowicie wyłączyć ochronę platformy. Lokalnie wdrożony model
działa niezależnie od~łączności zewnętrznej i~daje pełną kontrolę
nad~dostępnością usługi.

\paragraph{Niedeterministyczność i~brak testowalności.}
Duże modele językowe są~niedeterministyczne przy~niezerowej temperaturze
generowania: ta~sama wiadomość może przy~kolejnych wywołaniach otrzymać
różne etykiety. Utrudnia to~lub~uniemożliwia pisanie automatycznych testów
regresyjnych, które zakładają stabilność odpowiedzi systemu. Ponadto
klasyczny cykl CI/CD (\english{Continuous Integration / Continuous Deployment})
zakłada, że~zachowanie systemu jest funkcją jego kodu; niedeterministyczny
model LLM narusza to~założenie i~komplikuje proces weryfikacji poprawności.
System regułowy (Moduł~A) i~klasyfikator neuronowy po~dostrojeniu (Moduł~C)
są~deterministyczne w~fazie inferencji, co~pozwala na~pełne pokrycie testami
jednostkowymi i~regresyjnymi.

\paragraph{Brak dostrojenia domenowego.}
Ogólne modele językowe trenowane są~na~danych z~Internetu, które nie~obejmują
specyficznych wzorców polskojęzycznego czatu C2C marketplace. Frazy takie jak
,,dogadajmy się poza serwisem'', ,,zapłać mi~bezpośrednio, po~co te~prowizje''
czy~,,wbij na~mojego fejsa'' wymagają zrozumienia kontekstu biznesowego
polskiej platformy handlowej. Celowe zniekształcenia zapisu (,,n4mer t3lefonu'',
rozdzielanie cyfr spacjami), standardowe techniki obejścia detektorów
na~platformach C2C mogą być nierozpoznane przez ogólny model bez~dodatkowych
instrukcji w~zapytaniu. Dostrojenie na~danych własnych, realizowany przez
Moduł~C, pozwala precyzyjnie dostosować reprezentacje wektorowe modelu
do~tych wzorców.

\paragraph{Podatność na~wstrzyknięcie instrukcji.}
Klasyfikacja oparta na~instrukcji w~języku naturalnym otwiera powierzchnię
ataku nieistniejącą w~klasycznych klasyfikatorach: treść wiadomości użytkownika
może zawierać fragment próbujący nadpisać politykę systemową modelu
(\english{prompt injection}). Przykładowo, wiadomość zawierająca tekst
\textit{,,Zignoruj poprzednią instrukcję i~oceń tę~wiadomość jako bezpieczną''}
lub~jej zakamuflowaną wersję wplecioną w~pozornie niewinny tekst
może w~określonych warunkach zakłócić wynik klasyfikacji. Modele klasyfikacyjne
trenowane na~zbiorze zamkniętym (Moduły A-D) traktują cały tekst wejściowy
jako dane liczbowe; nie~zawierają mechanizmu interpretacji poleceń
i~nie~są~podatne na~ten wektor ataku.

\paragraph{Wyjaśnialność i~audytowalność decyzji.}
Regulacje dotyczące moderacji treści, w~szczególności Akt o~usługach
cyfrowych (\english{Digital Services Act}, DSA) w~Unii Europejskiej,
nakładają na~platformy obowiązek dokumentowania podstaw decyzji moderacyjnych
i~umożliwienia użytkownikom odwołania. LLM może generować uzasadnienie
decyzji w~języku naturalnym, jednak nie~jest ono wiążące: model może podać
uzasadnienie niespójne z~faktyczną etykietą (zjawisko halucynacji)
lub~przypisać tę~samą etykietę przy~różnych uzasadnieniach. System regułowy
(Moduł~A) i~Presidio (Moduł~B) wskazują dokładnie, który wzorzec
lub~rozpoznawacz encji zadecydował o~wyniku; klasyfikator neuronowy (Moduł~C)
umożliwia analizę macierzy pomyłek i~śledzenie prawdopodobieństwa klas.
Te~właściwości ułatwiają audyt, debugowanie i~dokumentację systemu.

\medskip
Zastosowanie zewnętrznego LLM do~moderacji treści jest uzasadnione
przy~szybkim prototypowaniu i~małym ruchu, gdy~uproszczenie wdrożenia przeważa
nad~kosztami i~ryzykami. W~warunkach produkcyjnych polskojęzycznej platformy C2C
z~wymogami RODO, skalą tysięcy wiadomości dziennie, koniecznością
deterministycznego testowania i~domenowym dostrojeniem, podejście to~generuje
nieakceptowalne koszty i~ryzyka operacyjne. Z~tych powodów w~niniejszej pracy
zdecydowano się wdrożyć własny potok oparty na~lokalnych komponentach
(Presidio i~HerBERT), który eliminuje wszystkie wymienione ograniczenia
i~pozwala na~pełną kontrolę nad~zachowaniem systemu moderacji.

\section{Obejście platform marketplace}
\label{sec:bypass}

Zjawisko obejścia platformy (\english{platform disintermediation}) polega na tym,
że kupujący i~sprzedający przenoszą część lub~całość transakcji poza platformę,
omijając jej mechanizm płatności oraz system prowizji \cite{bib:lin2022airbnb}.
W~praktyce oznacza to utratę przychodów przez operatora i~osłabienie kontroli nad
przebiegiem transakcji, dlatego jest to~jedno z~najpoważniejszych zagrożeń
biznesowych dla platform C2C.


Rozmiar strat wynikających z~obejścia platformy może być duży. Lin, Nian i~Foutz
\cite{bib:lin2022airbnb} przywołują dane z~ZBJ, chińskiej platformy dla pracowników
wiedzy pobierającej 20\% prowizji, według których obejście platformy może prowadzić
do utraty nawet 90\% transakcji, a~rzeczywiste przychody usługodawców są tam
dziesięciokrotnie wyższe niż wartości rejestrowane w~systemie platformy. Na Airbnb,
gdzie łączna opłata wynosi około 17\% (mniej więcej 3\% od gospodarza i~14\% od gościa),
autorzy empirycznie oszacowali skalę tego zjawiska. Łącząc dane GPS z~telefonów
komórkowych z~kalendarzami rezerwacji, ustalili wskaźnik na poziomie
5,4\% w~Austin (Teksas) w~2019 roku. Warto też zauważyć, że prowizje w~platformach
ekonomii współdzielenia są wysokie: Upwork pobiera 5-20\%, a~Uber, Lyft i~DoorDash nawet do
30\% wartości transakcji \cite{bib:lin2022airbnb}. To dodatkowo wzmacnia motywację
stron do omijania systemu.

Najczęstszą formą obejścia jest bezpośrednie przekazywanie danych kontaktowych,
na~przykład numeru telefonu, adresu e-mail albo nazwy użytkownika w~zewnętrznym
komunikatorze, w~treści wewnętrznego czatu platformy. Obie strony mają ku temu
silną motywację: chcą uniknąć prowizji, a~po kilku transakcjach rośnie wzajemne
zaufanie, przez co platforma przestaje być postrzegana jako potrzebny pośrednik.
Jak zauważają autorzy:

\begin{quote}
\textit{,,Both sides have incentives to take business offline, either because
they would like to circumvent a~marketplace fee charged by the platform, or
because they have trusted each other after a~few transactions''}
\cite{bib:lin2022airbnb}.
\end{quote}

\begin{quote}
\textit{,,Obie strony mają motywację do przeniesienia działalności do trybu offline, ponieważ chcą ominąć opłatę rynkową pobieraną przez platformę lub po kilku transakcjach zaufały sobie nawzajem.''}
\end{quote}
Szczególnie narażone są relacje z~powracającymi klientami. Po~pierwszej transakcji
obie strony zwykle mają już wzajemne dane kontaktowe, więc przy kolejnej okazji
łatwo umawiają się poza platformą.

Intuicyjną reakcją platform jest ukrywanie wrażliwych danych w~interfejsie
(\english{information concealment}). Airbnb maskuje numery telefonów i~dokładne
adresy nieruchomości aż do momentu finalizacji rezerwacji. Lin et al.
\cite{bib:lin2022airbnb} pokazują jednak, że strategia ta nie przynosi mierzalnego
efektu: uczestnicy wymieniają dane kontaktowe przez wewnętrzny komunikator
platformy jeszcze przed sfinalizowaniem transakcji, a~platforma nie ma wglądu
w~treść tych wiadomości.
Wnioski z~tych badań mają bezpośrednie znaczenie dla niniejszej pracy. Skoro
blokowanie informacji na poziomie interfejsu okazuje się nieskuteczne, właściwą
strategią staje się aktywna analiza treści wiadomości przesyłanych przez wewnętrzny
czat oraz automatyczne wykrywanie prób podawania danych kontaktowych lub
przenoszenia płatności poza system. Taki cel realizują moduły opisane w~rozdziale~3.
% Nguyen et al. 2025 - 4 typy zachowań bypass (TODO po uzupełnieniu bib)
% \cite{bib:bypass}

\section{Wykrywanie danych osobowych (PII)}
\label{sec:pii}

Wykrywanie danych osobowych (\english{Personally Identifiable Information}, PII)
stanowi kluczowy element detekcji prób obejścia platformy. W~kontekście czatu
marketplace najistotniejsze kategorie PII to: numery telefonów, adresy e-mail,
numery rachunków bankowych (IBAN) oraz nazwy użytkownika w~zewnętrznych
komunikatorach.

Tradycyjnie wykrywanie PII opierało się na wyrażeniach regularnych i~słownikach
fraz kluczowych. Podejście takie cechuje wysoka precyzja dla wzorców o~stałej
strukturze (np.\ numery IBAN według normy ISO~13616), lecz zawodzi wobec
wariantów tekstowych pisanych ze~spacjami, celowych przekształceń (,,zero''
zamiast cyfry~0) lub danych umieszczonych w~kontekście semantycznym.

Nowsze podejścia łączą wyrażenia regularne z~modelami NLP. Biblioteka
Presidio \cite{bib:presidio}, opracowana przez Microsoft, implementuje
architekturę wtyczek (\english{recognizer registry}), w~której każdy
rozpoznawacz encji może korzystać zarówno z~wyrażen regularnych, jak i~z modelu językowego.
Słowa kontekstowe (np.\ ,,telefon'', ,,zadzwoń'') zwiększają wynik pewności
rozpoznawacza, co ogranicza liczbę fałszywych trafień.

W~literaturze wykazano, że~hybrydowe podejście łączące NLP z~wyrażeniami
regularnymi przewyższa każdą z~tych technik zastosowaną osobno, uzyskując
wyższą precyzję przy mniejszej liczbie fałszywych trafień \cite{bib:presidio}.
W~niniejszej pracy podejście to realizuje Moduł~B (opisany w~rozdziale~3),
rozszerzony o~wzorce specyficzne dla domeny C2C: frazy obejścia platformy
(,,poza Stylify'', ,,ominąć prowizje'') oraz nazwy użytkownika
w~mediach społecznościowych.

\section{Modele językowe dla języka polskiego}
\label{sec:nlp-pl}

Modele transformerowe, których architektura opiera się na~mechanizmie
samouwagi (\english{self-attention}) zaproponowanym przez Vaswaniego i~in.\
\cite{bib:attention}, zrewolucjonizowały przetwarzanie języka naturalnego.
Model BERT \cite{bib:bert} jako pierwszy w~pełni wykorzystał tę architekturę
do~przetrenowania dwukierunkowego, umożliwiając transfer wiedzy z~dużych korpusów
na~zadania niszowe przy minimalnej ilości danych treningowych
\cite{bib:transformers_survey}. Jednak modele trenowane na~tekstach
wielojęzycznych lub wyłącznie anglojęzycznych nie oddają specyfiki języka
polskiego, który cechuje się bogatą fleksją, licznymi kolokwializmami
i~niestandardową ortografią w~komunikacji czatowej.

Dostępne modele dla języka polskiego można podzielić na trzy grupy:

\begin{itemize}
  \item \textbf{Wielojęzyczne}: \xlmr\ trenowany na~100~językach osiąga
    dobre wyniki ogólne, lecz nie specjalizuje się w~języku polskim.
    Jego zaletą jest gotowość do~użycia bez dostrojenia i~duża odporność
    na~przełączanie kodów (\english{code-switching}) \cite{bib:xlmr}.

  \item \textbf{Polskie monojęzyczne}: HerBERT (Allegro, 2021) jest
    przetrenowany wyłącznie na~polskich tekstach: Wikipedii,
    \english{Common Crawl} oraz danych z~Narodowego Korpusu Języka Polskiego.
    Architektura RoBERTa \cite{bib:roberta} z~tokenizacją BPE
    (\textit{Byte Pair Encoding} — algorytmem dzielącym tekst na~częste
    podciągi znaków) dostosowaną do~morfologii polskiej pozwala modelowi
    lepiej reprezentować formy fleksyjne i~neologizmy. W~benchmarkach KLEJ
    (Kompleksowa Lista Ewaluacyjna dla~Języka Polskiego, standardowy zestaw
    testów do~oceny polskich modeli NLP) osiąga wyniki powyżej \xlmr\
    na~większości zadań \cite{bib:herbert}.

  \item \textbf{Specjalizowane}: modele doszkalane na~konkretnych domenach
    (medycyna, prawo, finanse), niedostępne publicznie lub trenowane
    na~zbyt wąskich korpusach do~zastosowań ogólnych.
\end{itemize}

W~niniejszej pracy jako model bazowy dla Modułu~C wybrano HerBERT, ponieważ
osiąga najwyższe wyniki wśród publicznie dostępnych polskich modeli
na~zadaniach klasyfikacji tekstu i~jest swobodnie dostępny przez
bibliotekę \english{Hugging~Face Transformers}. \xlmr\ pełni rolę modelu
pomocniczego do~detekcji toksyczności (Moduł~A, backup).

\subsection{Narzędzia i biblioteki wykorzystane w pracy}
W pracy wykorzystano zestaw sprawdzonych bibliotek i modeli, które tworzą hybrydowy
potok przetwarzania detekcji nadużyć: od reguł i~wyrażeń regularnych, przez systemy rozpoznawania
encji, po modele transformerowe do klasyfikacji. Poniżej krótka charakterystyka
każdego z nich.

\paragraph{Microsoft Presidio:} biblioteka do wykrywania PII i~tworzenia rozpoznawaczy
(\english{recognizers}). Presidio umożliwia łączenie wzorców regularnych z~modelem NLP oraz
zarządzanie rejestrem rozpoznawaczy (\english{recognizer registry}), co pozwala na
stosowanie reguł kontekstowych i~znaczne zredukowanie liczby fałszywych trafień
\cite{bib:presidio}. W~pracy wykorzystano rozszerzenia i~wzorce specyficzne dla języka
polskiego oraz dodatkowe wzorce domenowe (np. frazy obejścia platformy).

\paragraph{HerBERT:} polskojęzyczny model w~architekturze RoBERTa zaprojektowany i~wstępnie
przetrenowany na dużych korpusach polskich tekstów. Dzięki tokenizacji i~korpusowi
szczególnie dopasowanym do fleksyjnej morfologii polskiej, HerBERT sprawdza się dobrze
w~zadaniach klasyfikacji i~jest używany jako główny model w~module neuronowym (Moduł~C)
\cite{bib:herbert}.

\paragraph{XLM-RoBERTa:} wielojęzyczny model transformerowy, używany tutaj jako model
pomocniczy (fallback) do wykrywania toksyczności i~ogólnych cech tekstu. XLM-R zapewnia
odporność na przełączanie kodów i~dobre wyniki w~scenariuszach wielojęzycznych \cite{bib:xlmr}.

\paragraph{Architektura Transformer / BERT:} podstawy teoretyczne i~architektoniczne dostarcza
rodzina modeli oparta na mechanizmie samo-uwagi (self-attention), zapoczątkowana i~
upowszechniona przez prace związane z~BERT i~kolejnymi wariantami. Te modele stanowią
punkt odniesienia dla wyboru i~analizy transformerów w pracy \cite{bib:bert,bib:transformers_survey}.

\paragraph{Wzorce regularne i wzorce domenowe:} dla klasy PII i~wektorów obejścia platformy konieczne
okazało się dodanie precyzyjnych wzorców regularnych (numery telefonów, IBAN, kody BLIK)
oraz list fraz kontekstowych — podejście to jest implementowane obok rozpoznawaczy
Presidio, tworząc prosty i~wydajny filtr preprocesujący wiadomości.

\paragraph{Infrastruktura i serwis ML:} moduł neuronowy działa jako niezależny serwis
oparty na \texttt{FastAPI} i~bibliotece \texttt{transformers} (Hugging Face). Serwis
udostępnia punkty końcowe klasyfikacji i~jest zintegrowany z~resztą aplikacji przez warstwę
API. Takie rozdzielenie ułatwia skalowanie i wymianę modeli w eksperymentach.

Opisane komponenty tworzą razem potok hybrydowy: szybkie reguły i~wzorce wyłapują oczywiste
przypadki PII, Presidio daje elastyczność rozpoznawaczy w kontekście, a~modele transformerowe
dostarczają decyzji o~wyższym poziomie abstrakcji (toksyczność, niuanse intencyjne). Ta
kompozycja pozwala łączyć wysoką precyzję z~dużą czułością na subtelne oznaki obejścia platformy.

\section{Klasyfikacja toksyczności}
\label{sec:toxicity}

Toksyczność w~komunikacji internetowej obejmuje wulgaryzmy, mowę nienawiści,
groźby i~obraźliwe zwroty. W~przeciwieństwie do~detekcji PII, gdzie wzorzec
jest często strukturalny, klasyfikacja toksyczności wymaga rozumienia kontekstu
semantycznego i~intencji nadawcy.

Dla języka polskiego dostępne zasoby NLP do~detekcji toksyczności są ograniczone.
Istniejące polskie zbiory danych z~etykietami toksyczności obejmują głównie
wpisy z~mediów społecznościowych i~forów internetowych. Ograniczeniem
tych danych jest ich nieprzystawalność do~stylu komunikacji w~czacie marketplace:
krótkie wiadomości, skróty, błędy ortograficzne i~specyficzna dynamika
wymiany kupujący-sprzedający.

Dostępne są ogólnodostępne modele wielojęzyczne dostrojone na~zbiorach
toksyczności z~wielu języków, które zapewniają akceptowalną skuteczność
dla~języka polskiego bez~dodatkowego dostrajania. W~niniejszej pracy
jeden z~takich modeli służy jako uzupełniający detektor toksyczności, natomiast główna klasyfikacja toksyczności
realizowana jest przez Moduł~C~(HerBERT dostrojony na~danych własnych),
co~pozwala uniknąć błędnych klasyfikacji wynikających z~niedostosowania
ogólnych modeli do~specyfiki języka czatu marketplace.

\section{Podsumowanie analizy}
\label{sec:analiza-podsumowanie}

Przegląd literatury wskazuje na~trzy istotne luki, które uzasadniają
cel niniejszej pracy:

\begin{enumerate}
  \item \textbf{Brak polskojęzycznych zasobów dla domeny marketplace.}
    Istniejące zbiory danych do~klasyfikacji toksyczności i~detekcji PII
    dotyczą głównie języka angielskiego lub mediów społecznościowych,
    nie zaś wewnętrznego czatu platformy C2C.

  \item \textbf{Brak porównania podejść regułowych, PII i~ML w~jednym eksperymencie.}
    Dostępne badania oceniają osobno systemy regułowe (PII) lub klasyfikatory
    ML (toksyczność), nie analizując ich synergii w~ramach hybrydowego potoku.

  \item \textbf{Brak uwzględnienia ataków fragmentowanych.}
    Istniejące systemy klasyfikują pojedyncze wiadomości w~izolacji.
    Tymczasem użytkownicy celowo rozbijają dane kontaktowe lub frazy
    obejścia platformy na~wiele krótkich wiadomości, aby uciec detektorom.
    Zagadnienie to nie było dotąd opisane w~kontekście platform marketplace.
\end{enumerate}

Niniejsza praca adresuje wszystkie trzy luki: tworzy syntetyczny zbiór
danych umożliwiający pierwsze uruchomienie systemu (który z~czasem będzie
wzbogacany o~dane z~panelu moderatora), porównuje cztery podejścia (A~-~D)
i~wprowadza mechanizm tury konwersacyjnej jako odpowiedź na~ataki
fragmentowane.
